\documentclass[a4paper,12pt]{article}
\usepackage{multirow}
\usepackage{enumitem}
\usepackage{adjustbox}
\usepackage{fontspec}
\usepackage{graphicx}
\usepackage[style="authoryear"]{biblatex}
\addbibresource{pustaka.bib}
\setmainfont{Times New Roman}
\graphicspath{ {./images/} }

%Setelan BibLatex
\AtBeginBibliography{%
  \renewcommand*{\finalnamedelim}{%
    \ifnumgreater{\value{liststop}}{2}{\finalandcomma}{}%
    \addspace dan \space}%
}
\renewbibmacro{in:}{}

%Setelan bagian
\renewcommand\thesection{\Alph{section}.}
\renewcommand\thesubsection{\thesection.\Alph{subsection}.}
\renewcommand\thesubsubsection{\thesection.\thesubsection. \Alph{susubbsection}.}

\newcommand{\bagian}[1]{\flushleft{\section{#1}}}
\newcommand{\subbagian}[1]{\flushleft{\subsection{#1}}}
\newcommand{\subsubbagian}[1]{\flushleft{\subsubsection{#1}}}

\newcommand{\addstretch}[1]{\addtolength{#1}{\fill}}
\newenvironment{autofull}
  {\newpage\flushbottom
   \addstretch{\baselineskip}
   \addstretch{\abovedisplayskip}
   \addstretch{\abovedisplayshortskip}
   \addstretch{\belowdisplayskip}
   \addstretch{\belowdisplayshortskip}
   \setlength{\parskip}{0pt}}
  {\newpage}

\begin{document}
\title{{\large PROPOSAL PRA SKRIPSI\\
  \Large{PENERAPAN MITM-PROXY UNTUK HTTPS FILTERING}}}
\author{Aldzikri Dwijayanto Prathama
    \\195410189}
\makeatletter
\begin{titlepage}
    \begin{center}
        {\huge \bfseries \@title}\\[14ex]
        \includegraphics[scale=.41]{logo}\\[4ex]
        {\large \bfseries Oleh:\\ \@author}\\[12ex]
        {\large \bfseries {PROGRAM STUDI INFORMATIKA\\
        PROGRAM SARJANA\\
        FAKULTAS TEKNOLOGI INFORMASI\\
        UNIVERSITAS TEKNOLOGI DIGITAL INDONESIA\\
        YOGYAKARTA\\
        2022}}
    \end{center}

\end{titlepage}
\makeatother

\newpage
\tableofcontents
\newpage

\center{\section*{Lembar Persetujuan}}
\vspace{5ex}
\begin{tabular}{rl}
  Judul    :& Penerapan MITM-Proxy Untuk HTTPS FILTERING\\[5ex]
  Nama     :& Aldzikri Dwijayanto Prathama\\[5ex]
  NIM      :& 195410189\\[5ex]
  Jurusan  :& Informatika\\[5ex]
  Semester :& VI(Enam)\\[5ex]
\end{tabular}
\begin{center}
  Telah memenuhi syarat dan disetujui untuk diseminarkan
  dihadapan dosen penguji seminar PROPOSAL PRA SKRIPSI.
\end{center}

\newpage
%\flushleft{\section{Latar Belakang}}
\bagian{Latar Belakang}
  \paragraph*{} Hampir seluruh perangkat saat ini terhubung
  dengan internet, seperti \textit{smartphone}, komputer,
  dan \textit{smart TV}. Meskipun tekonologi internet
  sangatlah bermanfaat, tetapi teknologi ini juga
  menimbulkan ancaman keamanan.

  \paragraph*{} Dari survey yang dilakukan, mayoritas
  developer tidak ragu melakukan \textit{personalized-ads}.
  Jenis iklan ini menggunakan data pengguna untuk memilihkan
  iklan yang dianggap relevan berdasarkan data pengguna. Hal
  ini tidak mengherankan karena jenis iklan ini menawarkan
  pendapatan yang lebih tinggi dibandingkan iklan biasa,
  karena peluang pengguna membuka iklan yang ditawarkan
  lebih besar(\cite{tahaei2021}). Penelitian lain juga
  menunujukkan bahwa 77\% aplikasi Android gratis
  menggunakan \textit{ad library}(\cite{he2018, jin2021}).

  \paragraph*{} Hal ini akan menjadi masalah ketika pihak
  yang mengumpulkan data tersebut mengalami kebocoran atau
  bahkan menggunakannya untuk hal-hal yang melanggar etika.

  \paragraph*{} Iklan-iklan yang ditampilkan sering kali
  sangat intrusif dan menggangu. Contohnya seperti iklan
  yang menutupi konten, iklan melalui \textit{pop-up}, atau
  iklan yang otomatis mengarahkan pengunjung ke \textit{tab}
  baru yang memuat iklan. Atau bahkan menampilkan iklan
  yang menjurus ke arah konten dewasa.

  \paragraph*{} Masalah lain selain iklan yang mengganggu,
  adalah masalah keamanan. Banyak layanan yang memasang
  iklan juga menggunakan layanan \textit{web analytics}
  untuk melakukan \textit{tracking} terhadap pengunjung atau
  penggunanya. Layanan web analytics inilah yang dapat
  menjadi celah keamanan. Contohnya pada tahun 2019, layanan
  Google Analytics digunakan untuk mengiklankan website
  phising(\cite{charlie2019}). Sedangkan di tahun 2020,
  peretas memasukkan kode berbahaya ke dalam website yang
  diretas, yang mana kode tersebut mengumpulkan informasi
  kredit pengguna dan mengirimkannya menggunakan Google
  Analytics, kemudian peretas akan mengakses data kredit
  yang dikumpulkan di akun Google Analytics
  miliknya(\cite{ravie2020}).

  \paragraph*{}Hal-hal tersebut bisa dicegah dengan
  menggunakan metode DNS \textit{ad-blocking}, dengan metode
  ini, kita bisa melakukan pemblokiran domain yang melakukan
  \textit{tracking}, domain yang digunakan sebagai
  tempat menyimpan kode JavaScript yang berbahaya, atau
  domain milik penyedia iklan. Tetapi metode
  ini hanya bisa melakukan pemblokiran berdasarkan nama
  domain, sehingga jika iklan atau kode JavaScript berada
  pada domain yang sama dengan website utama, maka kita
  harus mengizinkan iklan, dan kode JavaScript berjalan,
  atau memblokir domain website tersebut sehingga tidak bisa
  diakses.

  \paragraph*{}Metode lain adalah menggunakan
  \textit{add-on} peramban, seperti \textit{UBlock Origin}.
  Metode ini mampu melakukan \textit{blocking} terhadap
  iklan atau kode JavaScript yang berbahaya, walaupun
  keduanya berada pada domain yang sama dengan website
  utama. Sayangnya \textit{add-on} ini hanya bisa
  di-\textit{install} pada peramban tertentu, seperti
  Mozilla Firefox, dan Chrome Desktop.

  \paragraph*{}Berdasarkan permasalahan yang disebutkan
  di atas, tujuan saya adalah membangun sistem HTTPS
  filtering dengan menggunakan MITM-Proxy. Sistem ini
  bertujuan untuk menyaring konten yang tidak diinginkan
  agar tidak diakses.

  \paragraph*{}Filtering dengan menggunakan MITM-Proxy saya
  pilih karena tools ini merupakan proxy, sehingga bisa
  digunakan di banyak perangkat seperi metode DNS
  \textit{ad-blocking}, tetapi bisa di buat sedemikian rupa
  dengan menuliskan \textit{script} Python sehingga mampu
  melakukan filtering seperti \textit{add-on} pada browser,
  atau bahkan melebihinya.

  \paragraph*{}Berdasarkan penelitian yang sudah dilakukan
  sebelumnya, perbedaan pada penelitian ini adalah pada
  penggunaan MITM-Proxy. MITM-Proxy adalah HTTPS proxy
  interaktif yang gratis dan \textit{open source}. Tools ini
  bisa digunakan untuk \textit{debugging}, \textit{testing},
  dan \textit{penetration testing}. Software ini bisa
  digunakan sebagai alat untuk mencegat, inspeksi, dan
  memodifikasi lalu lintas web, seperti HTTP/1, HTTP/2,
  WebSockets, atau protokol lainnya yang dilindungi oleh
  SSL/TLS.

  \bagian{Rumusan Masalah}
  \paragraph*{}Berdasarkan latar belakang yang diuraikan di atas, dapat
  dirumuskan permasalahannya, yaitu "Bagaimana
  mengimplementasikan MITM-Proxy sebagai
  \textit{content-filtering}?"

  \printbibliography[title=Daftar Pustaka]

\end{document}
